\chapter{Strumenti matematici}
\label{cap:matematica}

\begin{inbreve}
Questo capitolo raccoglie gli strumenti usati in tutto il resto delle note. Sono il prodotto scalare e le norme,
le derivate parziali e il gradiente (la base della discesa del gradiente), le regole di derivazione per le
funzioni composte. Seguono i richiami di probabilità: gaussiana, valore atteso, varianza.
\end{inbreve}

\section{Vettori, prodotto scalare e norme}

Si usa la notazione: $x$ scalare, $\x$ vettore (colonna), $X$ matrice.

\begin{definizione}[Prodotto scalare]
Dati $\vv{a},\vv{b}\in\R^n$, il \termine{prodotto scalare} (\ingl{inner} o \ingl{dot product}) è
\begin{equation}
\vv{a}\cdot\vv{b} \;=\; a_1b_1 + a_2b_2 + \dots + a_nb_n \;=\; \sum_{i=1}^{n} a_i b_i \;=\; \norm{\vv{a}}\,\norm{\vv{b}}\cos\theta ,
\end{equation}
dove $\theta$ è l'angolo tra i due vettori. Notazioni equivalenti:
$\vv{a}\cdot\vv{b} = \vv{a}\T\vv{b} = \vv{a}^t\vv{b} = \langle\vv{a},\vv{b}\rangle$, e anche solo $\vv{a}\vv{b}$ quando il contesto è chiaro.
\end{definizione}

La forma con il coseno dà la lettura geometrica: il prodotto scalare misura \textbf{quanto due vettori puntano
nella stessa direzione}. È massimo quando sono paralleli e concordi ($\cos\theta = 1$), nullo quando sono
ortogonali ($\cos\theta=0$), negativo quando puntano in versi opposti. Nei modelli lineari, $\w\T\x$ misura
l'\emph{allineamento} dell'input $\x$ con il vettore dei pesi $\w$.

\begin{definizione}[Norme]
La \termine{norma euclidea} (modulo, lunghezza) di $\x$ è
\[
\norm{\x}_2 = \sqrt{\x\T\x} = \sqrt{\textstyle\sum_i x_i^2} = d(\x,\vv{0}),
\]
scritta spesso semplicemente $\norm{\x}$ o $\abs{\x}$. Altre norme utili sono la \termine{norma $L^1$},
$\norm{\x}_1 = \sum_i \abs{x_i}$, e la \termine{norma del massimo}, $\norm{\x}_\infty = \max_i\abs{x_i}$.
\end{definizione}

Più in generale, un \termine{prodotto interno} $\langle\cdot,\cdot\rangle$ su uno spazio vettoriale $V$ è una
forma bilineare simmetrica che associa a due vettori uno scalare, con le proprietà
\[
\langle v,w\rangle = \langle w,v\rangle, \qquad
\langle v+w,u\rangle = \langle v,u\rangle + \langle w,u\rangle, \qquad
\langle kv,w\rangle = k\,\langle v,w\rangle .
\]

\begin{teorema}[Disuguaglianza di Cauchy-Schwarz]
Per ogni $x,y\in V$ vale $\abs{\langle x,y\rangle} \le \norm{x}\cdot\norm{y}$, cioè, elevando al quadrato,
$\norm{x}^2\norm{y}^2 \ge \langle x,y\rangle^2$.
\end{teorema}

Sarà usata nella dimostrazione del teorema di convergenza del Perceptron (\rif{sec:convergenza-perceptron}).
Un'altra identità che serve spesso è lo sviluppo del quadrato di una norma:
\begin{equation}\label{eq:quadrato-norma}
\norm{\vv{a}+\vv{b}}^2 = (\vv{a}+\vv{b})\T(\vv{a}+\vv{b}) = \norm{\vv{a}}^2 + 2\,\vv{a}\T\vv{b} + \norm{\vv{b}}^2 .
\end{equation}

\begin{definizione}[Tensore]
Nella visione semplificata usata nel ML, un \termine{tensore} è un array multidimensionale di numeri disposti su
una griglia regolare. Il numero di assi è variabile: uno scalare ha 0 assi, un vettore 1, una matrice 2, e ad
esempio $X_{i,j,k}$ ha 3 indici. Nel ML il termine è soprattutto una comodità di notazione, utile per
descrivere il calcolo su GPU. In matematica i tensori hanno proprietà più ricche, legate ai cambi di base.
\end{definizione}

\section{Derivate parziali e gradiente}
\label{sec:gradiente}

La \termine{derivata parziale} generalizza la derivata alle funzioni di più variabili: è la derivata rispetto a
una sola variabile, tenendo fisse tutte le altre. Per $f(x_1,x_2,x_3)$ si calcolano
$\partial f/\partial x_1$, $\partial f/\partial x_2$, $\partial f/\partial x_3$.

\begin{definizione}[Gradiente]
Il \termine{gradiente} di $f:\R^n\to\R$ è il vettore delle sue derivate parziali:
\begin{equation}
\nabla f = \operatorname{grad} f = \left(\frac{\partial f}{\partial x_1},\dots,\frac{\partial f}{\partial x_n}\right)
= \sum_{i=1}^{n}\frac{\partial f}{\partial x_i}\,\vv{e}_i ,
\end{equation}
dove gli $\vv{e}_i$ sono i versori ortogonali lungo le direzioni coordinate.
\end{definizione}

\begin{idea}[Significato geometrico]
In un punto, il gradiente \textbf{punta nella direzione di massima pendenza} (massima crescita di $f$), e la sua
norma dice \textbf{quanto è ripida} la pendenza. Il suo opposto $-\nabla f$ indica la direzione in cui $f$
\textbf{decresce più rapidamente}. È l'idea alla base della \emph{discesa del gradiente}, che addestra quasi
tutti i modelli del corso minimizzando una funzione di errore.
\end{idea}

Il motivo si vede con il prodotto scalare. La variazione di $f$ spostandosi di poco lungo una direzione
unitaria $\vv{u}$ (la derivata direzionale) è
\[
D_{\vv{u}}f = \nabla f\cdot\vv{u} = \norm{\nabla f}\,\norm{\vv{u}}\cos\theta = \norm{\nabla f}\cos\theta ,
\]
con $\theta$ angolo tra $\vv{u}$ e $\nabla f$. Il valore è massimo per $\theta = 0$, cioè nella direzione di
$\nabla f$. Vale $-\norm{\nabla f}$, il minimo possibile, per $\theta=\pi$, cioè nella direzione di $-\nabla f$.
È nullo per $\theta=\pi/2$: lungo le \emph{curve di livello} $f$ non cambia, quindi il gradiente è
perpendicolare a esse (\figrif{fig:gradiente}).

\begin{figure}[t]
\centering
\begin{tikzpicture}
\begin{axis}[width=6.4cm, height=5.6cm, view={35}{30}, colormap={gr}{gray(0cm)=(0.92); gray(1cm)=(0.55)},
  axis lines=none, zmin=0, zmax=9]
\addplot3[surf, shader=faceted interp, faceted color=black!45, line width=0.1pt, domain=-2:2, y domain=-1.6:1.6, samples=17, samples y=15]
  {x^2+2*y^2};
\addplot3[only marks, mark=*, mark size=1.8pt, blu] coordinates {(0,0,0)};
\end{axis}
\end{tikzpicture}\hspace{0.6cm}
\begin{tikzpicture}[scale=1.15]
\foreach \c in {0.5,1,1.5,2,2.5}
  \draw[grigio] (0,0) ellipse ({\c} and {\c/1.414});
\draw[->, grigio!70] (-2.9,0) -- (2.9,0) node[right, black] {$x_1$};
\draw[->, grigio!70] (0,-2) -- (0,2.1) node[above, black] {$x_2$};
\fill[blu] (0,0) circle (1.6pt) node[above left, font=\footnotesize, fill=white, inner sep=1pt] {minimo};
% punto (1.2, 1.162) sulla curva c^2=4: x^2+2y^2=4
\coordinate (P) at (1.2,1.1662);
\fill (P) circle (1.5pt);
\draw[->, very thick, rosso] (P) -- ++(0.6,1.166) node[right, font=\small] {$\nabla f$};
\draw[->, very thick, verde] (P) -- ++(-0.6,-1.166) node[below right, font=\small] {$-\nabla f$};
\end{tikzpicture}
\caption{La funzione $f(x_1,x_2)=x_1^2+2x_2^2$. A sinistra la superficie, con il minimo (gradiente nullo); a
destra le curve di livello: in un punto $\nabla f$ (rosso) indica la crescita più rapida, $-\nabla f$ (verde) la
decrescita più rapida; entrambi sono perpendicolari alla curva di livello che passa per il punto.}
\label{fig:gradiente}
\end{figure}

\begin{definizione}[Punto stazionario]
Un punto in cui il gradiente è nullo, $\nabla f = \vv{0}$, si dice \termine{punto stazionario}. Può essere un
minimo locale, un massimo locale o un \termine{punto di sella}: il gradiente nullo, da solo, non basta a dire di
quale caso si tratta.
\end{definizione}

\subsection{Regole di derivazione}
\label{sec:regole-derivazione}

Il calcolo dei gradienti delle funzioni di errore usa sempre le stesse regole:
\begin{gather}
\frac{\partial k}{\partial w} = 0 \ (k \text{ costante}), \qquad \frac{\partial w}{\partial w} = 1, \qquad
\frac{\partial\, f(w)^2}{\partial w} = 2f(w)\frac{\partial f(w)}{\partial w}, \qquad
\frac{\partial}{\partial w}\sum_i f_i = \sum_i\frac{\partial f_i}{\partial w}, \\
\frac{\partial f}{\partial x} = \frac{\partial f}{\partial g}\cdot\frac{\partial g}{\partial x}
\ \text{(regola della catena)}, \qquad D[f(g(x))] = f'(g(x))\,g'(x), \qquad D[f\cdot g] = f'g + fg' .
\label{eq:catena}
\end{gather}
La regola della catena è la più importante: \textbf{scompone} il calcolo della derivata di una funzione composta
in fattori semplici, ciascuno calcolabile localmente. Su di essa si basa la backpropagation
(\rif{sec:backprop}).

\section{Probabilità}
\label{sec:probabilita}

\begin{definizione}[Valore atteso e varianza]
Per una variabile aleatoria discreta $Z$ che assume valori $z_i$ con probabilità $P(z_i)$:
\[
\bar Z = \E_P[Z] = \sum_i z_i\,P(z_i), \qquad \Var[Z] = \E\big[(Z-\bar Z)^2\big].
\]
Su un campione di $l$ valori si usano la media $\bar x = \frac1l\sum_p x_p$, la varianza
$\Var[x] = \frac1l\sum_p (x_p-\bar x)^2$ e la covarianza $\Cov[x,y] = \frac1l\sum_p (x_p-\bar x)(y_p-\bar y)$.
\end{definizione}

\begin{teorema}[Lemma della varianza]
\[
\Var[Z] = \E[Z^2] - \bar Z^2 \qquad\Longleftrightarrow\qquad \E[Z^2] = \bar Z^2 + \Var[Z].
\]
\end{teorema}
\begin{proof}
Si sviluppa il quadrato e si usa $\sum_i P(z_i) = 1$:
\begin{align*}
\Var[Z] &= \sum_i (z_i-\bar Z)^2P(z_i) = \sum_i z_i^2P(z_i) - 2\bar Z\sum_i z_iP(z_i) + \bar Z^2\sum_iP(z_i)\\
&= \E[Z^2] - 2\bar Z\cdot\bar Z + \bar Z^2 = \E[Z^2]-\bar Z^2. \qedhere
\end{align*}
\end{proof}

Il lemma è lo strumento chiave della decomposizione bias-varianza (\rif{sec:decomposizione}). Il valore atteso è
\textbf{lineare}: $\E[aX+bY] = a\E[X]+b\E[Y]$; se $X$ e $Y$ sono indipendenti vale anche $\E[XY]=\E[X]\,\E[Y]$.
Un campione si dice \termine{i.i.d.} (indipendente e identicamente distribuito) quando i suoi elementi sono
estratti indipendentemente l'uno dall'altro dalla stessa distribuzione.

\begin{definizione}[Distribuzione normale]
La \termine{distribuzione normale} (gaussiana) con media $\mu$ e varianza $\sigma^2$ ha densità
\begin{equation}
f(x) = \frac{1}{\sigma\sqrt{2\pi}}\;e^{-(x-\mu)^2/(2\sigma^2)} .
\end{equation}
Con $\mu=0$ e $\sigma^2=1$ si ottiene la \termine{normale standard}.
\end{definizione}

\begin{figure}[ht]
\centering
\begin{tikzpicture}
\begin{axis}[stile, width=9cm, height=4.6cm, domain=-3.6:3.6, ymin=0, ymax=0.45,
  xtick={-1.96,0,1,2,3,-1,-2,-3}, xticklabels={$-1{,}96$,$0$,$1$,$2$,$3$,$-1$,,$-3$}, ytick={0.1,0.2,0.3,0.4},
  yticklabels={$0{,}1$,$0{,}2$,$0{,}3$,$0{,}4$}, xlabel={$x$}, ylabel={$f(x)$}]
\addplot[name path=g, thick, inchiostro] {exp(-x^2/2)/sqrt(2*pi)};
\path[name path=ax] (axis cs:-3.6,0) -- (axis cs:3.6,0);
\addplot[pattern=north east lines, pattern color=rosso] fill between[of=g and ax, soft clip={domain=-3.6:-1.96}];
\node[font=\footnotesize, rosso, align=center] at (axis cs:-2.75,0.13) {area\\$2{,}5\%$};
\draw[densely dashed, grigio] (axis cs:-1.96,0) -- (axis cs:-1.96,0.0584);
\end{axis}
\end{tikzpicture}
\caption{Densità della normale standard ($\mu=0$, $\sigma=1$): l'area a sinistra di $-1{,}96$ vale il $2{,}5\%$.}
\label{fig:gaussiana}
\end{figure}

\begin{puntichiave}
\begin{itemize}
  \item $\vv{a}\T\vv{b} = \norm{\vv a}\norm{\vv b}\cos\theta$: il prodotto scalare misura l'allineamento.
  \item $\nabla f$ indica la crescita più rapida, $-\nabla f$ la decrescita più rapida; $\nabla f=\vv0$ individua
        i punti stazionari (minimi, massimi o selle).
  \item La regola della catena scompone le derivate di funzioni composte in fattori locali.
  \item $\E[Z^2] = \bar Z^2+\Var[Z]$.
\end{itemize}
\end{puntichiave}

\begin{domande}
\begin{enumerate}
  \item Che cos'è il gradiente e perché $-\nabla f$ indica la direzione di massima decrescita?
  \item Che cos'è un punto stazionario? È sempre un minimo?
\end{enumerate}
\end{domande}
